\documentclass[11pt,a4paper]{article}

\usepackage[utf8]{inputenc}
\usepackage[italian]{babel}
\usepackage{amsmath,amssymb,amsfonts,mathtools}
\usepackage{geometry}
\geometry{top=2.5cm, bottom=2.5cm, left=2.3cm, right=2.3cm}
\usepackage{booktabs}
\usepackage{tabularx}
\usepackage{hyperref}
\usepackage{listings}
\usepackage{xcolor}
\usepackage{microtype}
\usepackage{bm}

% Configurazione per il codice Python
\definecolor{codegreen}{rgb}{0,0.6,0}
\definecolor{codegray}{rgb}{0.5,0.5,0.5}
\definecolor{codepurple}{rgb}{0.58,0,0.82}
\definecolor{backcolour}{rgb}{0.97,0.97,0.97}

\lstdefinestyle{pythonstyle}{
    backgroundcolor=\color{backcolour},   
    commentstyle=\color{codegreen},
    keywordstyle=\color{blue}\bfseries,
    numberstyle=\tiny\color{codegray},
    stringstyle=\color{codepurple},
    basicstyle=\ttfamily\footnotesize,
    breakatwhitespace=false,         
    breaklines=true,                 
    captionpos=b,                    
    keepspaces=true,                 
    numbers=left,                    
    numbersep=5pt,                  
    showspaces=false,                
    showstringspaces=false,
    showtabs=false,                  
    tabsize=4,
    frame=single
}
\lstset{style=pythonstyle}

\title{\textbf{Architettura Computazionale e Visualizzazione delle Mappe di Confidenza $\Delta\chi^2$ 2D}\\\large Manuale Tecnico e Guida alle API del Modulo \texttt{chi\_squared\_2d\_map}}
\author{Laboratorio di Fisica Sperimentale e Computazionale}
\date{21 ottobre 2026}

\begin{document}

\maketitle

\begin{abstract}
Questo manuale fornisce la documentazione tecnica e operativa del modulo \texttt{chi\_squared\_2d\_map.py}, progettato per l'esplorazione e la rappresentazione grafica rigorosa delle superfici di costo bidimensionali scaturite da processi di fit ai minimi quadrati generalizzati (ODR Metodo~A). Vengono analizzati in dettaglio: l'algoritmo di profilazione vincolata con scansione a serpentina e \emph{warm-start} sequenziale di \texttt{Minuit}, l'interpolazione spline cubica per la ricostruzione continua della varietà, la topologia della griglia asimmetrica a quattro sub-riquadri via \texttt{GridSpec}, l'isolamento del box statistico di covarianza nel quadrante orfano e l'integrazione con i profili marginali 1D calcolati dall'algoritmo MINOS.
\end{abstract}

\tableofcontents
\newpage

% ──────────────────────────────────────────────────────────────────────────────
\section{Introduzione e Finalità del Modulo}
% ──────────────────────────────────────────────────────────────────────────────

Nelle analisi di fisica sperimentale, l'esame della sola matrice di covarianza $V = 2\mathbf{H}^{-1}$ (calcolata localmente nel minimo con \texttt{hesse}) fornisce una descrizione parabolica locale che ignora le asimmetrie e le correlazioni non lineari dello spazio dei parametri. Per visualizzare la reale geometria dell'incertezza congiunta di due parametri $(\theta_x, \theta_y)$, il modulo \texttt{chi\_squared\_2d\_map.py} implementa:
\begin{enumerate}
    \item \textbf{Scansione Profilata Reale (Profile Likelihood Scan)}: a ogni nodo della griglia $2\mathrm{D}$ i parametri $\theta_x$ e $\theta_y$ vengono fissati, mentre tutti i restanti $p-2$ parametri liberi del modello vengono rilassati e ri-ottimizzati tramite \texttt{migrad()}.
    \item \textbf{Sincronizzazione Geometrica 1D--2D}: proiezione sui bordi dei profili unidimensionali MINOS calcolati da \texttt{mnprofile}, verificando visivamente che le rette tangenti al contorno $1\sigma_{1\mathrm{D}}$ ($\Delta\chi^2 = 1.0$) coincidano con gli errori asimmetrici marginali.
    \item \textbf{Gestione Integrata dello Spazio Morto}: allocazione del box statistico con $\operatorname{Cov}(\theta_x, \theta_y)$ e $\rho_{xy}$ nel quadrante in alto a destra, eliminando interferenze con i dati.
    \item \textbf{Respiro Laterale della Colorbar}: posizionamento di un canale dedicato a destra con cuscinetto (\emph{buffer pad}), separando la barra di calibrazione continua dal profilo marginale verticale.
\end{enumerate}

% ──────────────────────────────────────────────────────────────────────────────
\section{Algoritmo di Profilazione e Ricostruzione della Superficie}
% ──────────────────────────────────────────────────────────────────────────────

\subsection{Campionamento a Serpentina e Warm-Start Sequenziale}

Dato un modello a $p$ parametri con minimo globale $F_{\mathrm{min}} = \chi^2(\hat{\boldsymbol{\theta}})$, si definisce un dominio di interesse bidimensionale centrato su $(\hat{\theta}_x, \hat{\theta}_y)$ di semi-ampiezza $N_\sigma$:
\begin{equation}
    x \in [\hat{\theta}_x - N_\sigma \sigma_x, \; \hat{\theta}_x + N_\sigma \sigma_x], \qquad y \in [\hat{\theta}_y - N_\sigma \sigma_y, \; \hat{\theta}_y + N_\sigma \sigma_y]
\end{equation}
dove $\sigma_x$ e $\sigma_y$ rappresentano le incertezze massime (o MINOS asimmetriche estese) estratte dall'oggetto \texttt{iminuit.Minuit}.

Il dominio viene discretizzato in una griglia regolare $K \times K$ (tipicamente $K \in [30, 45]$). La minimizzazione ripetuta a ogni nodo richiede una gestione accurata del punto iniziale per prevenire salti su minimi locali o fallimenti dell'algoritmo di discesa. Il modulo adotta una traiettoria a \textbf{serpentina alternata} (\emph{boustrophedon}):

\begin{lstlisting}[language=Python, caption={Scansione a serpentina con warm-start e fallback di ripristino.}]
for i in range(grid_size):
    # Alterna la direzione di scansione per massimizzare la vicinanza metrica
    j_indices = range(grid_size) if (i % 2 == 0) else reversed(range(grid_size))
    for j in j_indices:
        # 1. Warm-start: inizializza i p-2 parametri dallo stato ottimo del vicino
        for p in param_names:
            m.values[p] = current_cache[p]

        # 2. Fissazione della coppia bidimensionale
        m.fixed[param_x] = True
        m.fixed[param_y] = True
        m.values[param_x] = X_map[i, j]
        m.values[param_y] = Y_map[i, j]
        m.migrad()

        # 3. Fallback di sicurezza se la discesa diverge o non e' valida
        if not m.valid:
            for idx, p in enumerate(param_names):
                m.values[p] = orig_values[idx]
            m.values[param_x] = X_map[i, j]
            m.values[param_y] = Y_map[i, j]
            m.migrad()

        Z_profiled[i, j] = m.fval
        for p in param_names:
            current_cache[p] = m.values[p]
\end{lstlisting}

Il \emph{warm-start} riduce di oltre il $70\%$ le iterazioni interne di \texttt{migrad}, poiché il vettore dei parametri ottimi a coordinate $(i, j)$ costituisce un'eccellente approssimazione di partenza per il nodo adiacente $(i, j \pm 1)$.

\subsection{Regolarizzazione e Interpolazione Spline 2D}

La matrice campionata discretamente $Z_{ij} = \chi^2(x_i, y_j)$ viene riferita al minimo assoluto:
\begin{equation}
    \Delta Z_{ij} = Z_{ij} - F_{\mathrm{min}}
\end{equation}
Per tracciare curve di livello continue ed evitare seghettature poligonali nei contorni chiusi, la matrice $\Delta Z$ viene sovracampionata mediante interpolazione spline cubica (\texttt{scipy.ndimage.zoom}) con fattore di moltiplicazione $Z_f$ (default: 4):
\begin{lstlisting}[language=Python]
dz_smooth = ndimage.zoom(dz, zoom_factor, order=3)
\end{lstlisting}
La matrice continua risultante possiede dimensione $(K \cdot Z_f) \times (K \cdot Z_f)$, garantendo contorni ellittici fluidi e privi di artefatti numerici di discretizzazione.

\newpage

% ──────────────────────────────────────────────────────────────────────────────
\section{Topologia del Canvas e Geometria dei Subplot}
% ──────────────────────────────────────────────────────────────────────────────

\subsection{Partizione Asimmetrica con \texttt{gridspec.GridSpec}}

La disposizione simultanea della mappa 2D, dei due profili 1D e della barra colori richiede una suddivisione rigorosa dello spazio priva di automatismi coprenti. Con l'opzione standard \texttt{colorbar\_loc="right"}, la figura viene ripartita in una griglia di $2 \text{ righe} \times 4 \text{ colonne}$:

\begin{lstlisting}[language=Python, caption={Definizione della griglia proporzionale con colorbar a destra.}]
fig = plt.figure(figsize=(8.2, 7.2))
gs = gridspec.GridSpec(
    2, 4, 
    width_ratios=[3.5, 1.0, 0.1, 0.15], 
    height_ratios=[1.0, 3.5]
)
gs.update(wspace=spacing, hspace=spacing)

ax1  = fig.add_subplot(gs[1, 0])                 # Mappa principale 2D
ax_t = fig.add_subplot(gs[0, 0], sharex=ax1)     # Profilo 1D superiore
ax_r = fig.add_subplot(gs[1, 1], sharey=ax1)     # Profilo 1D laterale destro
cax  = fig.add_subplot(gs[1, 3])                 # Canale Colorbar dedicato
\end{lstlisting}

\begin{table}[htbp]
\centering
\small
\begin{tabularx}{\textwidth}{l c c X}
\toprule
\textbf{Subplot / Slot} & \textbf{Righe / Colonne} & \textbf{Ratio $(H \times W)$} & \textbf{Contenuto e Vincoli Fisici} \\
\midrule
\texttt{ax\_t} & $[0, 0]$ & $1.0 \times 3.5$ & Profilo $\Delta\chi^2$ parametro $x$; asse $x$ condiviso con \texttt{ax1}. \\
\emph{Quadrante Orfano} & $[0, 1]$ & $1.0 \times 1.0$ & Riquadro statistico: $\operatorname{Cov}(\theta_x, \theta_y)$, $\rho_{xy}$, errori asimmetrici. \\
\emph{Buffer Pad} & $[0, 2]$ e $[1, 2]$ & $1.0 \times 0.1$ & Intervallo vuoto di distacco per isolare la colorbar. \\
\texttt{ax1} & $[1, 0]$ & $3.5 \times 3.5$ & Mappa di densità $\Delta\chi^2$, livelli di contorno e punto ottimo $(+)$. \\
\texttt{ax\_r} & $[1, 1]$ & $3.5 \times 1.0$ & Profilo $\Delta\chi^2$ parametro $y$; asse $y$ condiviso con \texttt{ax1}. \\
\texttt{cax} & $[1, 3]$ & $3.5 \times 0.15$ & Scala cromatica graduata continua con tacche sui livelli $\sigma_{2\mathrm{D}}$. \\
\bottomrule
\end{tabularx}
\caption{Mappatura delle celle della griglia \texttt{GridSpec} in configurazione laterale.}
\end{table}

\subsection{Sincronizzazione degli Assi e Soppressione Sovrapposizioni}
La coerenza visiva esige l'azzeramento dello spazio tra la mappa e i profili (\texttt{wspace=0.01}, \texttt{hspace=0.01}):
\begin{itemize}
    \item L'asse orizzontale di \texttt{ax\_t} condivide i limiti con \texttt{ax1} (\texttt{sharex=ax1}), ma le etichette numeriche inferiori vengono soppresse:
    \begin{lstlisting}[language=Python]
ax_t.tick_params(axis='x', bottom=True, labelbottom=False)
    \end{lstlisting}
    \item L'asse verticale di \texttt{ax\_r} condivide i limiti con \texttt{ax1} (\texttt{sharey=ax1}), sopprimendo le etichette numeriche a sinistra:
    \begin{lstlisting}[language=Python]
ax_r.tick_params(axis='y', left=True, labelleft=False)
    \end{lstlisting}
    \item \textbf{Soppressione del Tick Zero Sovrapposto}: per evitare che la cifra \texttt{0} sull'asse $\Delta\chi^2$ di \texttt{ax\_r} si sovrapponga visivamente all'ultimo tick numerico dell'asse $x$ di \texttt{ax1}, è possibile applicare un formattatore di potatura:
    \begin{lstlisting}[language=Python]
ax_r.xaxis.set_major_formatter(
    ticker.FuncFormatter(lambda v, _: "" if abs(v) < 1e-3 else f"{v:.0f}")
)
    \end{lstlisting}
\end{itemize}

\newpage

% ──────────────────────────────────────────────────────────────────────────────
\section{Livelli di Confidenza, Contorni e Segnaletica}
% ──────────────────────────────────────────────────────────────────────────────

\subsection{Significato Statistico dei Valori $\Delta\chi^2$}

Nella rappresentazione congiunta di due parametri, la variabile $\Delta\chi^2(\theta_x, \theta_y)$ segue asintoticamente una distribuzione chi-quadro a $\nu = 2$ gradi di libertà:
\begin{equation}
    P(\Delta\chi^2 \le \Delta\chi^2_0) = \int_0^{\Delta\chi^2_0} \frac{1}{2} e^{-z/2}\,dz = 1 - e^{-\Delta\chi^2_0 / 2}
\end{equation}
Al contrario, per un profilo marginale monodimensionale (proiettato su un singolo asse rilassando l'altro), i gradi di libertà sono $\nu = 1$:
\begin{equation}
    P(\Delta\chi^2 \le \Delta\chi^2_0) = \operatorname{erf}\left(\sqrt{\frac{\Delta\chi^2_0}{2}}\right)
\end{equation}

\begin{table}[htbp]
\centering
\small
\begin{tabularx}{\textwidth}{l c c c X}
\toprule
\textbf{Livello Tracciato} & \textbf{Gradi di Libertà} & $\Delta\chi^2$ & \textbf{Copertura} & \textbf{Stile Grafico Adottato} \\
\midrule
$1\sigma_{\text{1D}}$ (Marginale) & $\nu = 1$ & $1.00$ & $68.27\%$ & Tratto-punto nero (\texttt{'-.'}, spessore 1.1 pt) \\
$1\sigma_{\text{2D}}$ (Congiunto) & $\nu = 2$ & $2.30$ & $68.27\%$ & Continuo bianco (\texttt{'-'}, spessore 1.6 pt) \\
$2\sigma_{\text{2D}}$ (Congiunto) & $\nu = 2$ & $6.18$ & $95.45\%$ & Tratteggiato bianco (\texttt{'--'}, spessore 1.2 pt) \\
$3\sigma_{\text{2D}}$ (Congiunto) & $\nu = 2$ & $11.83$ & $99.73\%$ & Punteggiato bianco (\texttt{':'}, spessore 0.9 pt) \\
\bottomrule
\end{tabularx}
\caption{Livelli di confidenza standard e stili geometrici implementati nel modulo.}
\end{table}

\subsection{Verifica Geometrica delle Tangenti MINOS}

Sul piano $2\mathrm{D}$, il contorno ad anello $1\sigma_{\text{1D}}$ ($\Delta\chi^2 = 1.00$) non racchiude il $68.27\%$ della probabilità congiunta (che richiederebbe $\Delta\chi^2 = 2.30$), bensì possiede una proprietà geometrica fondamentale:
\begin{quote}
\emph{Le rette tangenti orizzontali e verticali al contorno $\Delta\chi^2 = 1.00$ individuano esattamente gli intervalli di incertezza marginale a $1\sigma$ calcolati dal solutore MINOS per i singoli parametri.}
\end{quote}
Il modulo traccia tali tangenti mediante linee rosse tratteggiate (\texttt{ls='--', lw=1.2}), che si prolungano simultaneamente sia nella mappa centrale sia nei subplot 1D marginali, evidenziando visivamente la concordanza tra la scansione $2\mathrm{D}$ e i profili parabolici/asimmetrici $1\mathrm{D}$.

\subsection{Leggibilità del Testo con Effetto Contorno (\emph{Halo Stroke})}

Per garantire che le etichette numeriche dei livelli di contorno (es. $1\sigma_{\text{2D}}, 2\sigma_{\text{2D}}$) rimangano perfettamente visibili sia su sfondi scuri che su regioni a saturazione luminosa (come accade con le mappe \texttt{viridis\_r} o \texttt{plasma}), il testo viene corredato da un effetto contorno (\emph{path stroke}):

\begin{lstlisting}[language=Python, caption={Applicazione del filtro tipografico halo con \texttt{patheffects}.}]
halo = [patheffects.withStroke(linewidth=2.5, foreground="white")]
clbl_2d = ax1.clabel(cs_2d, inline=True, fmt=fmt_clabel, colors="black")
plt.setp(clbl_2d, path_effects=halo)

clbl_1d = ax1.clabel(cs_1d, inline=True, fmt=fmt_clabel, colors="black")
plt.setp(clbl_1d, path_effects=halo)
\end{lstlisting}

\newpage

% ──────────────────────────────────────────────────────────────────────────────
\section{Firma Completa delle API: \texttt{chi2\_map\_pipeline}}
% ──────────────────────────────────────────────────────────────────────────────

\begin{lstlisting}[language=Python, caption={Firma della funzione principale \texttt{chi2\_map\_pipeline}.}]
def chi2_map_pipeline(
    risultati: Any,
    param_x: str,
    param_y: str,
    param_info: Optional[Dict[str, Dict[str, str]]] = None,
    n_sigma: float = 3.5,
    grid_size: int = 35,
    zoom_factor: int = 4,
    show_clabels: bool = True,
    manual_clabels: Union[bool, List[Tuple[float, float]]] = False,
    clabel_fontsize: float = typo_params["legendsize"],
    stats_box_loc: str = "inside_map",  # 'inside_map', 'top', o 'custom'
    stats_box_pos: Optional[Tuple[float, float]] = None,
    show_stats_box: bool = True,
    spacing: float = 0.01,
    cmap: str = "viridis_r",
    save_path: Optional[str] = None,
    show: bool = True,
    colorbar_loc: str = "right",        # 'right', 'bottom', oppure 'none'
) -> Tuple[plt.Figure, Tuple[plt.Axes, plt.Axes, plt.Axes]]:
\end{lstlisting}

\subsection{Dizionario dei Parametri di Configurazione}

\begin{table}[htbp]
\centering
\small
\begin{tabularx}{\textwidth}{l p{2.8cm} X}
\toprule
\textbf{Parametro} & \textbf{Tipo / Default} & \textbf{Descrizione e Funzione Operativa} \\
\midrule
\texttt{risultati} & \texttt{RisultatiFit} & Contenitore dell'output del fit contenente l'istanza \texttt{minuit\_obj}. \\
\texttt{param\_x}, \texttt{param\_y} & \texttt{str} & Nomi delle due variabili fisiche da scansionare (es. \texttt{"m"}, \texttt{"q"}). \\
\texttt{param\_info} & \texttt{dict} / \texttt{None} & Metadati di visualizzazione: simbolo \LaTeX\ (\texttt{'label'}) e unità di misura (\texttt{'unit'}). \\
\texttt{n\_sigma} & \texttt{float = 3.5} & Semi-estensione dell'intervallo di scansione espressa in deviazioni standard. \\
\texttt{grid\_size} & \texttt{int = 35} & Risoluzione della griglia nodale $K \times K$ prima della regolarizzazione. \\
\texttt{zoom\_factor} & \texttt{int = 4} & Fattore di interpolazione spline cubica ($K \cdot Z_f$ punti finali). \\
\texttt{show\_clabels} & \texttt{bool = True} & Abilita la scrittura a video delle etichette dei contorni ($1\sigma_{\text{2D}}$, ecc.). \\
\texttt{manual\_clabels} & \texttt{bool / list} & Coordinate manuali delle etichette per evitare sovrapposizioni su curve ripide. \\
\texttt{stats\_box\_loc} & \texttt{str = "inside\_map"} & Posizione del box di correlazione: \texttt{'inside\_map'}, \texttt{'top'}, \texttt{'custom'}. \\
\texttt{stats\_box\_pos} & \texttt{tuple} / \texttt{None} & Coordinate $(x, y)$ in frazione di figura per la modalità \texttt{'custom'}. \\
\texttt{colorbar\_loc} & \texttt{str = "right"} & Posizione della scala colori: \texttt{'right'}, \texttt{'bottom'} o disabilitata. \\
\texttt{spacing} & \texttt{float = 0.01} & Separazione orizzontale e verticale tra i riquadri (\texttt{wspace, hspace}). \\
\texttt{cmap} & \texttt{str = "viridis\_r"} & Mappa cromatica continua di Matplotlib per i valori del $\Delta\chi^2$. \\
\texttt{save\_path} & \texttt{str / Path / None} & Percorso su disco per l'esportazione automatica in formato vettoriale PDF/SVG. \\
\bottomrule
\end{tabularx}
\caption{Parametri di configurazione del generatore di mappe profilate.}
\end{table}

\newpage

% ──────────────────────────────────────────────────────────────────────────────
\section{Integrazione con la Pipeline di Fit ed Esempio Completo}
% ──────────────────────────────────────────────────────────────────────────────

Lo script seguente illustra il ciclo completo di laboratorio: generazione dati con errori bivariati, stima non distorta con \texttt{FitPipeline} (ODR Metodo~A) e produzione della mappa con box statistico nel quadrante in alto a destra (\texttt{stats\_box\_loc="custom"}).

\begin{lstlisting}[language=Python, caption={Flusso end-to-end: regressione ortogonale e rendering della mappa 2D.}]
import numpy as np
from pathlib import Path
from chi_squared_fit import FitPipeline
from chi_squared_2d_map import chi2_map_pipeline

# 1. Definizione della funzione teorica (Legge di Ohm con offset)
def modello_lineare(x, m, q):
    return m * x + q

# 2. Generazione del campione sintetico con errori su entrambe le variabili
rng = np.random.default_rng(123)
m_vero, q_vero = 2.45, 1.10
n_punti = 35

x_pulito = np.linspace(1.0, 10.0, n_punti)
y_pulito = modello_lineare(x_pulito, m_vero, q_vero)

err_x = np.full(n_punti, 0.20)
err_y = np.full(n_punti, 0.45)

x_misurato = x_pulito + rng.normal(0, err_x, n_punti)
y_misurato = y_pulito + rng.normal(0, err_y, n_punti)

# 3. Esecuzione FitPipeline con Metodo A (Varianza Congelata e MINOS)
pipeline = FitPipeline(
    model=modello_lineare,
    param_names=["m", "q"],
    p0=[1.0, 0.0],
    run_minos=True,
    verbose=True,
)

risultati_fit = pipeline.fit(
    x=x_misurato,
    y=y_misurato,
    err_x=err_x,
    err_y=err_y,
    odr_lin_Newton=True,  # Attiva il Metodo A non distorto
)

# 4. Configurazione dei metadati dimensionali
info_parametri = {
    "m": {"label": r"$m$", "unit": r"\mathrm{V/s}"},
    "q": {"label": r"$q$", "unit": r"\mathrm{V}"},
}

# 5. Generazione della Mappa 2D ad alta risoluzione
save_folder = Path("examples")
save_folder.mkdir(parents=True, exist_ok=True)

fig, (ax_main, ax_top, ax_right) = chi2_map_pipeline(
    risultati=risultati_fit,
    param_x="m",
    param_y="q",
    param_info=info_parametri,
    n_sigma=3.2,
    grid_size=32,
    zoom_factor=4,
    stats_box_loc="custom",
    stats_box_pos=(0.725, 0.85),  # Quadrante in alto a destra
    show_clabels=True,
    manual_clabels=False,
    colorbar_loc="right",
    cmap="viridis_r",
    spacing=0.0,
    save_path=save_folder / "mappa_chi2_profilata.pdf",
    show=True,
)
\end{lstlisting}

\end{document}